% Modernised reading — fonds Alexandre Grothendieck, shelfmark 143, pages 1-33.
%
% This is an INTERPRETATION, not a transcription. It re-reads the pages in
% today's notation and vocabulary, states the hypotheses the manuscript leaves
% implicit, and drops the critical apparatus so that the mathematics reads
% straight through. Everything the brackets and underlines carried moves into a
% footnote. It is held to being mathematically correct as it stands; where the
% manuscript is loose or mistaken the true statement appears and the footnote
% says what the page has. In any dispute of substance the transcription
% governs: whoever cites Grothendieck cites batch-01.fr.tex and batch-02.fr.tex.
%
% Scope: the folder is thirty-three pages in two batches (1-20, 21-33), both
% transcribed, read whole before any of this was written, and written as ONE
% file for the whole shelfmark. Page 33 is a computer listing with nothing of
% his; page 32 is written over another listing.
%
% The spine. One text, his pp. 1-16 (archivists' 1-30 with the unnumbered
% versos), plus two loose sheets (31, 32):
%   - pp. 1-10 (his §§ 1-6): the « paradigm » data of (X, D, J) over Q-bar —
%     the group G of absolute automorphisms, the extension Sigma of G by pi_1,
%     germs of liftings attached to points, the Kummer structure at the
%     components of D, the cusps X~'(I) with Phi, the marked points X~'(J)
%     with the partial splittings Psi; recapitulated I / II / III, (1)-(9);
%   - p. 11: the « Programme des Gleichnis » (a), (b), and the NB;
%   - pp. 12-13: RemB, descent to a number field and bitorsors B_gamma;
%   - pp. 14-18 (§ 7): the fundamental groupoid Pi_1(X'; J) rebuilt from
%     X~'(J); « Reformulation »: the equivalence alpha / beta between
%     extensions-with-a-Sigma-set and groupoids-with-a-universal-cover;
%   - pp. 18-26 (§ 8): the cusp datum (T-script, T, Phi) in the same terms,
%     the action groupoid <T-script, pi>, T-groupoids, the pseudo-partition
%     condition, the functor E |-> T-script^0 x_pi E; « Arrêtons »;
%   - pp. 26-27 (his « (9) »): modular topoi of genus 0;
%   - pp. 28-31 (« Une conjecture arithmético-géométrique », also numbered 9):
%     pi_1(M_{g,nu}) -> Out_lac°(pi^_{g,nu}), diagram (4), sequence (5), NB,
%     two questions;
%   - p. 32: the forgetful fibration M_{g,nu+1} -> M_{g,nu} and its profinite
%     counterpart.
%
% Notation fixed for the whole folder. X is the Q-bar-scheme (the pages write
% X-bar, then a round X); X' = X - supp D; pi_1 is the étale fundamental group,
% X~ a universal (pro-)covering. G-script = absolute automorphisms, G = the
% Q-bar-linear ones; Sigma the extension. T = T(Q-bar) = lim mu_n(Q-bar) =
% Z^(1), chi the cyclotomic character. Groups act on the left; the right action
% of pi_1 on a left pi_1-set is x.lambda = lambda^-1 x (needed for contracted
% products, as in 141). In the abstract part, 1 -> pi -> Sigma -> G -> 1, Z
% (the page's curly 3) for the Sigma-set modelling X~'(J), T-script for the one
% modelling X~'(I); the orbit set T-script/pi is written I here (the pages
% write J). Autext is written Out, as in 145; Out_lac with the page's
% exponents ° (each inertia class fixed) and °° (and multiplier 1). T_{g,nu} is
% his « groupe de Teichmüller », today's pure mapping class group.
%
% Substantive departures from the pages, each footnoted where it occurs:
%   - p. 1: G-script -> Gamma needs X connected (the margin says a condition is
%     needed); the convention making it a homomorphism is ours; the lifting
%     determines the K-form without « réduit », and a lifting comes from a
%     K-form only if continuous and, for effectivity, X quasi-projective;
%   - p. 4: « xi is known from its germ » is proved here (Kummer theory of a
%     semi-abelian variety, no divisible points over a number field);
%   - p. 6, (19): the target Out(pi_1; T, (phi_i)) is defined here;
%   - p. 8, (25): X~(j) is X~'(j); Hom(G_j, Sigma') is the set of sections;
%   - p. 10, (9): the germ Gamma' -> Sigma' depends on a point (it is (15));
%   - p. 14: Isom(X~'[j], X~'[j']) is paths j' -> j under Yoneda; harmless;
%   - p. 15: X~'(I) in the bracket is X~'(J);
%   - p. 17: j |-> Z_j is covariant (the margin writes Pi° -> Ens); Pi~ is the
%     indiscrete groupoid on Z, a discrete fibration over Pi;
%   - pp. 19-21: the orbit set T-script/pi is I, not J; the stabiliser pi_t
%     centralises Phi(t)(T) (proved), so « T commutatif » is forced by
%     Phi_t : T ~ pi_t; Sigma acts on T through chi, not trivially;
%   - p. 26: M_{0,nu}, nu <= 2, are classifying stacks of positive-dimensional
%     groups, not DM; p. 27: pi_1 = 0 is the geometric pi_1; the braid group
%     is P_{nu-1}/centre, the pure mapping class group of the sphere;
%   - p. 28: the brace under (***) belongs to Out°°, not to pi_1;
%   - p. 30: « bijectif » fails for (g, nu) = (0, 4); p. 31: Out(T_g) = Z/2 for
%     g >= 3 (Ivanov, McCarthy), not 1.
%
% Announced and never established in the folder: the careful treatment of dim
% X > 1 (margins of pp. 10, 11); the explicit passage to finite étale covers
% (« bouclage » / « fixage de tresses », p. 11); the « variantes profinies »
% of the reformulation (p. 18); an interpretation of the T-structure (« ceci
% ne donne pas d'interprétation … Arrêtons », p. 26); the list of « cas
% typiques » promised on p. 11; the conjecture of pp. 28-30 and the two
% questions of pp. 30-31.
%
% Pass: Opus 5.5 (claude-opus-5-5), 2026-10-10 — first pass, unchecked against the pages by a human.
\documentclass[11pt,a4paper]{article}
\input{../preamble/grothendieck}

\folder{143}
\pages{1}{33}
\dating{[vers 1980-1981]}
\watermark{Édition de démonstration\\interprétation personnelle de l'œuvre}
\foldertitle{Paradigmes des courbes algébriques (version provisoire) : notes manuscrites (s.d.) — lecture modernisée du dossier entier}

\begin{document}

\section*{Résumé}

\begin{resume}
Ces trente-deux pages manuscrites dressent un inventaire : tout ce que le
groupe de Galois des nombres algébriques « voit » d'une courbe algébrique,
quand on ne regarde que ses lacets.

Une courbe complexe privée de quelques points est une surface trouée. Les
trajets fermés qu'on peut y tracer à partir d'un point forment un groupe, le
\emph{groupe fondamental}, et ce groupe, pris avec toutes ses limites
(le groupe « profini »), connaît tous les revêtements finis de la courbe. Si la
courbe est définie par des équations à coefficients algébriques, les symétries
des nombres algébriques — le groupe de Galois $\Gamma$, qui échange par exemple
$\sqrt{2}$ et $-\sqrt{2}$ — ne déplacent pas la surface comme une figure qu'on
dessine, mais elles agissent sur ses lacets. La question des feuillets est :
qu'est-ce qui, dans cette action, est structure, et peut-on tout décrire
avec des groupes seuls ?

Grothendieck en fait la liste. Il y a l'extension qui mêle les lacets et les
symétries ; il y a, pour chaque trou, le petit lacet qui en fait le tour, que
la symétrie galoisienne envoie sur un petit lacet autour d'un trou, mais
parcouru un nombre de fois qui mesure son action sur les racines de l'unité
(la « structure kummérienne ») ; il y a les points de la courbe, dont chacun
fournit une manière de relever les symétries dans le groupe des lacets. Il
récapitule le tout en une liste de neuf éléments de structure, puis écrit un
programme en deux lignes, sous le titre allemand de « Gleichnis » —
paraboles, ou paradigmes : décrire complètement $\Gamma$ en termes de groupes,
et décrire complètement toute la structure inventoriée.

Viennent ensuite des pages plus abstraites, où l'on voit l'auteur chercher la
bonne formulation. Une famille de points de base, avec les lacets qui les
joignent, forme un \emph{groupoïde} — un groupe à plusieurs objets — et il
montre que se donner une extension de groupes avec un ensemble sur lequel elle
agit, ou se donner un groupoïde avec son revêtement universel, c'est la même
chose. Il tente d'y faire entrer les petits lacets autour des trous, obtient
une formulation correcte, la juge « pas bien éclairante », en essaie une autre,
et s'arrête : « Arrêtons… ». Le geste est caractéristique : ce n'est pas
l'exactitude d'un énoncé qui le satisfait, mais le moment où il devient
évident.

Les dernières pages changent d'échelle. Au lieu d'une courbe, on prend
l'espace de \emph{toutes} les courbes d'un genre donné avec des points
marqués, l'espace des modules. Ses lacets agissent sur les lacets de chaque
courbe, et Grothendieck conjecture que cette action donne exactement toutes
les symétries des lacets qui respectent les trous. Pour la sphère privée de
trois points, cela dirait que le groupe de Galois $\Gamma$ est le groupe de
ces symétries — c'est la question qu'étudie en détail le dossier voisin 145.
La conjecture n'a pas été démontrée ; elle a été raffinée, après lui, en ce
qu'on appelle la géométrie anabélienne et le groupe de
Grothendieck--Teichmüller. Une question qu'il se pose en passant — le groupe
des lacets de l'espace des modules se plonge-t-il dans les symétries des
lacets de la courbe ? — est aujourd'hui connue sous le nom de problème des
sous-groupes de congruence.

Pour aller plus loin : groupe fondamental arithmétique, action galoisienne
extérieure, conjecture des sections, groupoïde fondamental, bitorseur, champ
des modules des courbes, groupe modulaire de Teichmüller, suite exacte de
Birman.

\keywords{anabelian geometry, arithmetic fundamental group, outer Galois action, section conjecture, cuspidal inertia, cyclotomic character, Galois descent, fundamental groupoid, gauge groupoid, covering of groupoids, banded gerbe, moduli stack of curves, Deligne–Mumford stack, mapping class group, Dehn–Nielsen–Baer theorem, Birman exact sequence, congruence subgroup problem}
\end{resume}

\pagerange{1}{33}
\section*{Le fil du dossier, et les conventions}

\subsection*{Les stations}

Le dossier est une seule rédaction, paginée de sa main de 1 à 16 (sur les
rectos, les versos courant sans numéro), suivie de deux feuillets isolés. Il
numérote ses paragraphes de ① à ⑧, puis deux fois ⑨.

\begin{itemize}
\item \textbf{Les données d'un paradigme (pages 1 à 10, ses pages 1 à 5)} :
le groupe $\mathcal{G}$ des automorphismes absolus d'un schéma $X$ sur
$\overline{\mathbb{Q}}$, l'extension $\Sigma$ de $\mathcal{G}$ par
$\pi_1(X)$, le germe de relèvement attaché à un point, la structure
kummérienne attachée à un diviseur $D$, les pointes $\widetilde X'(I)$ et
l'application $\Phi$, les points marqués $\widetilde X'(J)$ et les scindages
$\Psi$ ; puis la récapitulation en trois groupes I, II, III et neuf éléments
de structure.
\item \textbf{Le programme (page 11)} : décrire $\Gamma$ et toute cette
structure « en termes de théorie des groupes ».
\item \textbf{Descente et bitorseurs (pages 12 et 13)} : comment la donnée
précédente contient l'action de $\Gamma'$ sur les revêtements.
\item \textbf{Le groupoïde fondamental (pages 14 à 18)} : $\Pi_1(X'; J)$
reconstruit à partir de $\widetilde X'(J)$, puis la « Reformulation » :
extensions munies d'un ensemble, ou groupoïdes munis d'un revêtement
universel.
\item \textbf{Les pointes en termes de groupoïdes (pages 18 à 26)} :
plusieurs essais pour dire la structure $\Phi$ dans ce langage ; le dernier
est abandonné.
\item \textbf{Topos modulaires de genre $0$ (pages 26 et 27).}
\item \textbf{Une conjecture arithmético-géométrique (pages 28 à 31)} :
l'action de $\pi_1(M_{g,\nu})$ sur le groupe fondamental de la courbe
universelle serait un isomorphisme sur le groupe des automorphismes
extérieurs « à lacets ».
\item \textbf{La fibration d'oubli d'un point (page 32)}, deux diagrammes
sur un listing d'ordinateur. La page 33 est un listing qui ne porte rien de
lui.
\end{itemize}

Le dossier n'établit pas de lien explicite entre la première moitié et la
seconde. Il y en a un, que nous signalons comme nôtre là où il sert : les
données des pages 1 à 13, écrites pour une courbe $X'$, sont induites, pour
la courbe $X_\xi$ correspondant à un point $\xi$ de l'espace des modules, par
la fibration de la page 32 ; et le groupe $\pi_1(N_{g,\nu})$ de la page 30 est
l'extension $\Sigma$ des pages 1 et 2, écrite pour l'espace des modules.

\subsection*{Les conventions}

\emph{Schémas et groupes fondamentaux.} $X$ désigne un schéma de type fini sur
$\overline{\mathbb{Q}}$, connexe et non vide\footnote{La page 1 écrit
$\overline{X}$, puis passe dès le ② à une lettre ronde ; on écrit $X$ partout.
Dans les dernières pages, $X$ redevient une courbe, celle qui correspond à un
point de l'espace des modules : ce sont les mêmes objets.}. $\pi_1$ est le
groupe fondamental étale, profini, et $\widetilde X \to X$ un revêtement
universel, c'est-à-dire un pro-revêtement étale galoisien de groupe
$\pi_1(X) = \mathrm{Aut}(\widetilde X/X)$. Le groupe $\pi_1$ agit à gauche ; un
$\pi_1$-ensemble à gauche $E$ devient à droite par $x\cdot\lambda =
\lambda^{-1}x$, et $A \wedge^{\pi_1} E$ est le produit contracté, quotient de
$A \times E$ par $(a\lambda, e) \sim (a, \lambda e)$ — les conventions du
dossier 141.

\emph{Galois.} $\Gamma = \mathrm{Gal}(\overline{\mathbb{Q}}/\mathbb{Q})$ ;
$\Gamma'$ désigne un sous-groupe ouvert, celui d'un corps de nombres $K$.
$T = T(\overline{\mathbb{Q}}) = \varprojlim_n \mu_n(\overline{\mathbb{Q}})$, qu'on
écrit aujourd'hui $\widehat{\mathbb{Z}}(1)$, et $\chi : \Gamma \to
\widehat{\mathbb{Z}}^{*} = \mathrm{Aut}(T)$ le caractère cyclotomique.

\emph{Automorphismes extérieurs.} On écrit $\mathrm{Out}$ pour son
« $\mathrm{Autext}$ », comme dans la lecture du dossier 145, et
$\mathrm{Out}_{\mathrm{lac}}$ pour les automorphismes extérieurs « à lacets »
(on dit aujourd'hui : qui préservent l'inertie) — ceux qui envoient chaque
classe de conjugaison de sous-groupes d'inertie des trous sur une telle
classe, avec un même multiplicateur dans $\widehat{\mathbb{Z}}^{*}$.
L'exposant ${}^{\circ}$ de la page 28 impose de fixer chaque classe, ce que le
dossier 145 note ${}^{!}$ ; l'exposant ${}^{\circ\circ}$ impose en outre le
multiplicateur $1$.

\emph{Le modèle abstrait.} À partir de la page 16, $1 \to \pi \to \Sigma \to
G \to 1$ est une extension abstraite de groupes ; $\mathfrak{Z}$ est un
$\Sigma$-ensemble qui modèle $\widetilde X'(J)$, $\mathcal{T}$ un
$\Sigma$-ensemble qui modèle $\widetilde X'(I)$. Les pages appellent $J$ à la
fois $\mathfrak{Z}/\pi$ et $\mathcal{T}/\pi$ ; le second correspond, dans le
dictionnaire de la page 8, à l'ensemble $I$ des composantes de $D$, et on
l'écrit $I$\footnote{Pages 19, 20 et 21. Le nom $J$ n'y est pas faux en
soi — c'est un nom neuf pour un ensemble neuf —, mais il ferait lire les
pointes comme des points marqués.}.

\emph{Espaces de modules.} $M_{g,\nu}$ est le champ (le « topos modulaire »)
des courbes lisses propres de genre $g$ munies de $\nu$ sections ordonnées
deux à deux disjointes, sur $\mathbb{Q}$ ; $\overline{M}_{g,\nu}$ est
$M_{g,\nu} \otimes \overline{\mathbb{Q}}$ — la barre est l'extension des
scalaires, non la compactification. $T_{g,\nu}$ est ce qu'il appelle le
« groupe de Teichmüller », le groupe modulaire (mapping class group) pur de la
surface de genre $g$ à $\nu$ trous, et $\widehat T_{g,\nu}$ son complété
profini ; $\widehat\pi_{g,\nu}$ est le complété profini du groupe fondamental
de cette surface.

\pagerange{1}{1}
\section*{Le groupe des automorphismes absolus (page 1)}

Soit $X$ comme ci-dessus et $\mathcal{G} = \mathcal{G}_X$ le groupe des
automorphismes de $X$ en tant que schéma \emph{absolu}, c'est-à-dire sans
exiger qu'ils respectent la structure de $\overline{\mathbb{Q}}$-schéma. Un tel
automorphisme $u$ agit sur l'anneau $\mathcal{O}(X)$ des fonctions, et
préserve le sous-anneau des éléments annulés par un polynôme séparable à
coefficients rationnels ; comme $X$ est connexe, ce sous-anneau est
exactement $\overline{\mathbb{Q}}$ (une telle fonction prend ses valeurs parmi
les racines distinctes d'un polynôme, et chaque valeur découpe une partie
ouverte et fermée). D'où une suite exacte
\[
(1)\qquad 1 \longrightarrow G \longrightarrow \mathcal{G} \longrightarrow
\Gamma, \qquad G = \mathrm{Aut}(X/\overline{\mathbb{Q}}),
\]
où $u$ s'envoie sur l'élément $\gamma$ tel que $u^{*}$ agisse sur
$\overline{\mathbb{Q}}$ par $\gamma^{-1}$\footnote{La note marginale de droite
porte « $\overline{\mathbb{Q}} \hookrightarrow \Gamma(X, \mathcal{O}_X)$ … il
faut … » : une condition est nécessaire pour que la flèche existe. Celle que
nous donnons — $X$ connexe — est de cette édition. L'inverse est une
convention, la page n'en fixe pas : $u \mapsto u^{*}|_{\overline{\mathbb{Q}}}$
renverse l'ordre des produits, et son inverse est un homomorphisme. Avec elle,
$\mathcal{G}$ agit sur $X(\overline{\mathbb{Q}})$ par $\xi \mapsto u \circ \xi
\circ \mathrm{Spec}(\gamma)$, et pour $X = X_K \otimes_K \overline{\mathbb{Q}}$
on retrouve l'action galoisienne usuelle sur les coordonnées.}.

\emph{L'image est d'indice fini.} Les sous-extensions $K$ de
$\overline{\mathbb{Q}}/\mathbb{Q}$ correspondent aux sous-groupes fermés
$\Gamma_K$, d'indice fini si et seulement si $K$ est fini sur $\mathbb{Q}$.
Comme $X$ est de type fini, ses équations n'ont qu'un nombre fini de
coefficients : $X$ provient d'un $K$-schéma $X_K$ pour un corps de nombres
$K$, et $\gamma \mapsto \mathrm{id}_{X_K} \times \mathrm{Spec}(\gamma^{-1})$
est un relèvement
\[
(4)\qquad \Gamma_K \longrightarrow \mathcal{G}
\]
de l'inclusion $\Gamma_K \subset \Gamma$. L'image de $\mathcal{G} \to \Gamma$
contient donc le sous-groupe ouvert $\Gamma_K$ : elle est ouverte, d'indice
fini\footnote{La marge gauche pose la question (« d'indice fini ? … vrai ») ;
le texte la tranche, et il a raison.}.

Le relèvement (4) \emph{détermine} la $K$-structure : les fonctions de $X_K$
sont les fonctions de $X$ invariantes par $\Gamma_K$, et cela sans hypothèse
de réduction\footnote{La page dit « qui le définit si $\overline{X}$ est
réduit et loc. de t.f. », et cercle d'un « ? » le mot « réduit » : son doute
est fondé, la descente galoisienne n'en a pas besoin.}. Inversement, un
relèvement $\Gamma_K \to \mathcal{G}$ ne vient d'une $K$-forme que s'il est
\emph{continu} (chaque fonction a une orbite finie), et, même alors, la donnée
de descente n'est effective que sous une hypothèse comme
« $X$ quasi-projectif » : c'est le « grain de sel d'effectivité … de
projectivité » de la note marginale\footnote{La note est au crayon, en
oblique, et ne se lit qu'en partie ; « projectivité » est une lecture
incertaine. Pour un schéma quasi-projectif, toute donnée de descente galoisienne
continue est effective ; en général, elle ne l'est pas.}.

\pagerange{1}{4}
\section*{L'extension $\Sigma$, les points et leurs germes (pages 1 à 4)}

\pagerange{1}{2}
\subsection*{L'extension $\Sigma_X$}

Soit $\widetilde X \to X$ un revêtement universel, et
\[
(5)\qquad \Sigma_X = \Sigma(X, \widetilde X) = \mathrm{Aut}(\widetilde X /\!/ X)
\]
le groupe des automorphismes du schéma absolu $\widetilde X$ qui recouvrent un
élément de $\mathcal{G}$\footnote{Le signe entre $\widetilde X$ et $X$ est une
barre oblique surchargée, rendue $/\!/$ par la transcription ; la définition
ci-dessus est ce que la suite exacte (6) impose. La lettre $\Sigma$ hésite sur
la page entre un E rond et un sigma.}. Si $u \in \mathcal{G}$, le composé
$\widetilde X \to X \xrightarrow{u} X$ est encore un revêtement universel, donc
$u$ se relève, et les relèvements forment un torseur sous $\pi_1(X) =
\mathrm{Aut}(\widetilde X/X)$ :
\[
(6)\qquad 1 \longrightarrow \pi_1(X) \longrightarrow \Sigma_X \longrightarrow
\mathcal{G} \longrightarrow 1 .
\]
La conjugaison dans $\Sigma_X$ fait agir $\mathcal{G}$ sur $\pi_1(X)$ par
automorphismes extérieurs. Au-dessus de $\Gamma_K$ (par (4)), $\Sigma_X$ n'est
autre que le groupe fondamental du $K$-schéma $X_K$, dans la suite exacte
fondamentale $1 \to \pi_1(X) \to \pi_1(X_K) \to \Gamma_K \to 1$
(c'est ce que reprend la page 12). C'est la même construction que le groupe
$\pi_1(X, \Gamma, a)$ de la page 14 du dossier 141 et que l'extension
$\mathcal{E} = \pi_1(U, \Gamma; P)$ « à la Malgoire » des pages 15 et 16 du
dossier 145, en version étale et pour le groupe $\mathcal{G}$ tout
entier\footnote{Le rapprochement est de cette édition. On dirait aujourd'hui
groupe fondamental du champ quotient $[X/\mathcal{G}]$, ou groupe fondamental
équivariant.}.

\emph{Un point.} Soit $\xi \in X(\overline{\mathbb{Q}})$ et
$\mathcal{G}(\xi)$ son stabilisateur dans $\mathcal{G}$. Si $G$ est fini,
$\mathcal{G}$ est profini et $\mathcal{G}(\xi)$ est un sous-groupe ouvert :
$\mathcal{G}$ est alors extension d'un sous-groupe ouvert de $\Gamma$ par un
groupe fini, et $\mathcal{G}(\xi)$ contient l'image de $\Gamma_{K'}$ dès que
$\xi$ est rationnel sur $K'$\footnote{L'hypothèse « $X$ de type fini sur
$\overline{\mathbb{Q}}$ » est ajoutée par lui en interligne ; elle est
nécessaire.}. Choisissons $\widetilde X$ pointé, c'est-à-dire un point
$\tilde\xi$ de $\widetilde X$ au-dessus de $\xi$. Pour $u \in \mathcal{G}(\xi)$,
les relèvements de $u$ permutent simplement transitivement la fibre de
$\widetilde X$ en $\xi$ ; un seul fixe $\tilde\xi$. D'où un scindage partiel
de (6), qui est un homomorphisme :
\[
(8)\qquad \mathcal{G}(\xi) \longrightarrow \Sigma_X, \qquad u \longmapsto
\text{l'unique relèvement de } u \text{ fixant } \tilde\xi .
\]
Changer $\tilde\xi$ le compose avec un automorphisme intérieur provenant de
$\pi_1(X)$ : sans point choisi, (8) n'est défini qu'à cette conjugaison
près\footnote{La phrase biffée qui commence au bas de la page 2 et se
poursuit, barrée, en haut de la page 3 annonçait déjà la condition de la page
4 : « s'envoie injectivement dans un groupe algébrique commutatif de rang
unipotent nul (ext. d'une VA par un tore) ».}.

\pagerange{3}{4}
\subsection*{Le germe de relèvement, et ce qu'il dit du point (pages 3 et 4)}

Deux $K$-formes de $X$ deviennent égales sur une extension finie $K'$ de
$K$ : l'identité de $X$ est définie sur un corps de nombres. Les relèvements
(4) qu'elles définissent coïncident donc sur $\Gamma_{K'}$. Ce qui est
canonique est ainsi un \emph{germe} de relèvement
\[
(9)\qquad \Gamma' \longrightarrow \mathcal{G}, \qquad \Gamma' \subset \Gamma
\text{ ouvert},
\]
défini modulo la relation « coïncider sur un sous-groupe ouvert plus petit ».
Pour $\Gamma'$ assez petit, il se factorise par $\mathcal{G}(\xi)$, et le
composé avec (8) donne un germe
\[
(10)\qquad \varphi_\xi : \Gamma' \longrightarrow \Sigma_X, \qquad
\varphi_\xi \in \varinjlim_{\Gamma' \subset \Gamma \text{ ouvert}}
\mathrm{Hom}(\Gamma', \Sigma_X),
\]
relevant l'inclusion de $\Gamma'$, défini à conjugaison près par
$\pi_1(X)$ si le point $\tilde\xi$ n'est pas fixé. « Scholie : … encore une
fois que c'est le germe seulement de $\varphi_\xi$ qui est
défini »\footnote{Le mot qui suit « Scholie : » est illisible. La formule (11)
de la page 4 porte, après $\mathrm{Hom}(\Gamma', \Sigma_X)$, un argument
noirci et barré puis un groupe $(\pi_1(X), \Sigma_X)$ précédé d'un signe
illisible ; on n'en tient pas compte.}. En termes d'aujourd'hui :
si $\xi$ est rationnel sur $K'$, $\varphi_\xi$ est la \emph{section} de la suite
exacte fondamentale $1 \to \pi_1(X) \to \pi_1(X_{K'}) \to \Gamma_{K'} \to 1$
associée au point rationnel $\xi$, bien définie à conjugaison près.

La page ajoute entre crochets : \emph{si $X$ s'envoie injectivement dans une
extension d'une variété abélienne par un tore, $\xi$ est connu quand on connaît
son germe} — et l'image de $\varphi_\xi$ dans un $H^1$ convenable suffit. C'est
vrai, et voici pourquoi. Soit $\iota : X \hookrightarrow A$ un tel plongement
dans une variété semi-abélienne, $\xi, \xi'$ deux points rationnels sur $K'$
dont les germes coïncident à conjugaison près sur un $\Gamma_{K''}$. Poussées
dans $\pi_1(A_{K''})$, les deux sections diffèrent d'une classe de
$H^1(K'', \widehat{T}(A))$, où $\widehat{T}(A) = \pi_1(A)$ est le module de
Tate total, et cette classe est l'image de $\iota(\xi) - \iota(\xi')$ par
l'application de Kummer $A(K'') \to H^1(K'', \widehat{T}(A))$. Elle est
nulle ; or le noyau de l'application de Kummer est la partie infiniment
divisible de $A(K'')$, et elle est triviale sur un corps de nombres : l'image
d'un tel élément dans la variété abélienne quotient est nulle par le théorème
de Mordell--Weil, et l'on se ramène au tore, où le groupe multiplicatif d'un
corps de nombres n'a pas d'élément infiniment divisible autre que $1$. Donc
$\iota(\xi) = \iota(\xi')$, et $\xi = \xi'$\footnote{La démonstration est de
cette édition ; la page affirme, entre crochets. C'est l'énoncé
d'\emph{injectivité} de ce qu'on appelle aujourd'hui la conjecture des
sections, que Grothendieck formulera dans sa lettre à Faltings de 1983 : pour
une courbe propre hyperbolique sur un corps de nombres, l'application des points
rationnels vers les classes de sections serait bijective. L'injectivité est
connue par l'argument ci-dessus ; la surjectivité est ouverte.}.

\pagerange{4}{7}
\section*{Un diviseur : structure kummérienne et pointes (pages 4 à 7)}

\pagerange{4}{6}
\subsection*{L'extension $\Sigma_{X'}$ et les homomorphismes de Kummer}

Soit $D = \sum_{i \in I} D_i$ un diviseur effectif réduit sur $X$, $I$
l'ensemble de ses composantes irréductibles, et supposons $X$ régulier aux
points génériques $\eta_i$ des $D_i$. Posons
\[
(13)\qquad X' = X \setminus \operatorname{supp} D,
\]
et soit $\mathcal{G}_{X'}$ le sous-groupe de $\mathcal{G}_X$ qui préserve
$D$\footnote{Il le note aussi $\mathcal{G}_X(I)$, « abus de notation ». Un
mot souligné sous l'indice se lit, mal, « ouvert » : $\mathcal{G}_{X'}$ est
en effet ouvert dans $\mathcal{G}_X$ quand celui-ci est profini, puisqu'il
contient un $\Gamma_{K}$ pour lequel $D$ est défini sur $K$.}. Avec un
revêtement universel $\widetilde X'$ de $X'$, la construction précédente donne
\[
(14)\qquad 1 \longrightarrow \pi_1(X') \longrightarrow \Sigma_{X'}
\longrightarrow \mathcal{G}_{X'} \longrightarrow 1,
\]
et, pour $\xi \in X'(\overline{\mathbb{Q}})$, un germe $\varphi_\xi : \Gamma'
\to \Sigma_{X'}$ (15). Le groupe $\mathcal{G}_{X'}$ permute les composantes de
$D$ :
\[
(16)\qquad \mathcal{G}_{X'} \longrightarrow \mathfrak{S}_I .
\]

\emph{La structure kummérienne.} L'anneau local $\mathcal{O}_{X,\eta_i}$ est un
anneau de valuation discrète, de corps résiduel de caractéristique $0$. Son
hensélisé strict a pour corps des fractions un corps dont toutes les
extensions finies sont kummériennes, de groupe de Galois $\varprojlim_n
\mu_n = T$ : c'est la situation de la page 1 du dossier 141, corps des
fractions d'un anneau de valuation discrète hensélien à corps résiduel
algébriquement clos de caractéristique nulle\footnote{Les racines de l'unité
de ce corps résiduel sont celles de $\overline{\mathbb{Q}}$, qui y est
algébriquement fermé : le groupe d'inertie est donc canoniquement
$T(\overline{\mathbb{Q}})$, ce qui justifie l'écriture de la page.}. Le choix
d'un point de $\widetilde X'$ « au-dessus de $\eta_i$ » donne un
homomorphisme de ce groupe d'inertie dans $\pi_1(X')$ ; sans ce choix, on a un
homomorphisme extérieur
\[
(17)\qquad \varphi_i : T \longrightarrow \pi_1(X'), \qquad \varphi_i \in
\mathrm{Homext}(T, \pi_1(X')) = \mathrm{Hom}(T, \pi_1(X'))/\mathrm{Int}(\pi_1(X')),
\]
la classe du lacet qui fait le tour de $D_i$\footnote{Une note au crayon de
plusieurs lignes, à gauche de (17), où se lit « (17) », est hachurée et ne se
lit pas.}. Le groupe $\mathcal{G}_{X'}$ agit sur $T$ à travers $\Gamma$ et le
caractère cyclotomique :
\[
(18)\qquad \mathcal{G}_{X'} \longrightarrow \Gamma \xrightarrow{\ \chi\ }
\widehat{\mathbb{Z}}^{*} = \mathrm{Aut}(T).
\]
L'action extérieure de $\mathcal{G}_{X'}$ sur $\pi_1(X')$ est compatible avec
la famille $(\varphi_i)$ — la « structure kummérienne » — compte tenu des
actions (16) et (18) : pour $u \in \mathcal{G}_{X'}$, on a $u_* \circ \varphi_i
\equiv \varphi_{u(i)} \circ \chi(u)^{\pm 1}$ modulo automorphismes intérieurs,
le signe dépendant de la convention de (1). On peut donc écrire
\[
(19)\qquad \mathcal{G}_{X'} \longrightarrow \mathrm{Out}\bigl(\pi_1(X'); T,
(\varphi_i)\bigr),
\]
où le but est le groupe des triplets $(a, s, c) \in \mathrm{Out}(\pi_1(X'))
\times \mathfrak{S}_I \times \mathrm{Aut}(T)$ tels que $a \circ \varphi_i \equiv
\varphi_{s(i)} \circ c$ pour tout $i$\footnote{La définition du but est de
cette édition ; la page écrit $\mathrm{Autext}(\pi_1(X'), T(\overline{\mathbb{Q}}),
(\varphi_i))$, le deuxième argument surchargé et de lecture incertaine. Pour une
courbe, c'est le groupe $\mathrm{Out}_{\mathrm{lac}}$ des conventions, avec
$c$ pour multiplicateur ; c'est lui qui revient page 28.}.

\pagerange{6}{7}
\subsection*{Les pointes, et l'application $\Phi$ (pages 6 et 7)}

Soit $\widehat{\widetilde X}{}' \to X$ le normalisé de $X$ dans le
pro-revêtement $\widetilde X'$ de $X'$, morphisme profini, et
$\widetilde X'(I)$ l'image inverse de l'ensemble des points maximaux de $D$,
c'est-à-dire des $\eta_i$\footnote{« génériques » est biffé, « maximaux »
écrit en interligne.}. Pour une courbe, ce sont les \emph{pointes} du
revêtement universel. On a une application
\[
(20)\qquad \widetilde X'(I) \longrightarrow I,
\]
sur laquelle $\Sigma_{X'}$ agit, compatiblement avec (16) : un élément de
$\Sigma_{X'}$ préserve $D$ et s'étend donc au normalisé. La fibre
$\widetilde X'(i)$ en $i$ est un espace homogène sous $\pi_1(X')$, puisque le
groupe d'un revêtement galoisien permute transitivement les points au-dessus
d'un point donné\footnote{Note à droite de (21), qui s'interrompt : « C'est
un espace homogène sous $\pi_1(X')$. On … ».}. À un point $\tilde\eta$ de
$\widetilde X'(i)$, associons l'homomorphisme d'inertie de $T$ dans son
stabilisateur, vu dans $\pi_1(X')$ :
\[
(22)\qquad \Phi_i : \widetilde X'(i) \longrightarrow \mathrm{Hom}(T, \pi_1(X')).
\]
Il est compatible avec $\pi_1(X')$ des deux côtés — $\Phi_i(\lambda\tilde\eta) =
\mathrm{int}(\lambda) \circ \Phi_i(\tilde\eta)$ — et précise l'homomorphisme
extérieur $\varphi_i$ de (17), qui en est la classe. Mieux, le système
\[
(23)\qquad \Phi = (\Phi_i)_{i \in I} : \widetilde X'(I) \longrightarrow
\mathrm{Hom}(T, \pi_1(X'))
\]
est compatible avec $\Sigma_{X'}$ des deux côtés, $\Sigma_{X'}$ agissant sur le
but par conjugaison sur $\pi_1(X')$ et par (18) sur $T$ :
$\Phi(\sigma\tilde\eta) = \mathrm{int}(\sigma) \circ \Phi(\tilde\eta) \circ
\chi(\sigma)^{\mp 1}$\footnote{Une note de plusieurs lignes en oblique, en
marge gauche de la page 7, est presque entièrement illisible : « ici il
faudrait … dire … $X$ … $\widetilde X'(I)$ … $\Phi$ … ».}.

L'homomorphisme $\Phi(\tilde\eta)$ n'est pas injectif en général —
il est nul pour $X' = \mathbb{A}^1$ —, et son image n'est pas tout le
stabilisateur de $\tilde\eta$ dès que $\dim X > 1$, le corps résiduel de
$\eta_i$ ayant alors lui-même un groupe de Galois. Pour une courbe, il en va
autrement, et c'est ce qui servira page 20.

\pagerange{7}{8}
\section*{Des points marqués (pages 7 et 8)}

Soit enfin, « par dessus le marché », une partie $J \subset
X'(\overline{\mathbb{Q}})$, $\widetilde X'(J)$ son image inverse dans
$\widetilde X'$, et $\mathcal{G}' = \mathcal{G}_{X'}(I, J)$ le sous-groupe de
$\mathcal{G}_{X'}$ qui préserve $J$. Il agit sur $J$, sur $\widetilde X'(J)$ et
sur
\[
(24)\qquad \widetilde X'(J) \longrightarrow J .
\]
Notons $\Sigma'$ l'image inverse de $\mathcal{G}'$ dans $\Sigma_{X'}$ et, pour
$j \in J$, $\mathcal{G}'_j$ le stabilisateur de $j$ et $\Sigma'_j$ son image
inverse. Le scindage (8) de la page 2 donne une famille d'applications
\[
(25)\qquad \Psi = (\Psi_j)_{j \in J}, \qquad \Psi_j : \widetilde X'(j)
\longrightarrow \mathrm{Splitt}(\Sigma'_j/\mathcal{G}'_j),
\]
où $\mathrm{Splitt}(\Sigma'_j/\mathcal{G}'_j)$ est l'ensemble des sections
homomorphes de $\Sigma'_j \to \mathcal{G}'_j$, et $\Psi_j(\tilde\jmath)(\sigma)$
l'unique relèvement de $\sigma$ qui fixe $\tilde\jmath$\footnote{La page 8 écrit
$\widetilde{X}(j)$, sans prime, et le but
$\mathrm{Hom}(\mathcal{G}_j, \Sigma')$ ; la page 10 le corrige en
$\mathrm{Splitt}(\Sigma'_j/\mathcal{G}'_j)$, qu'on adopte. La suite de la page
8, barrée de deux longs traits, esquissait une variante à valeurs dans
$\mathrm{Hom}(\mathcal{G}', \Sigma')/\mathrm{Int}(\pi_1, \Sigma')$ ; elle revient,
barrée, comme (8) de la page 10.}. Elle est compatible avec $\Sigma'_j$, qui
agit sur le but par
\[
({}^{g}u)(\sigma) = g\, u(g_1^{-1}\sigma g_1)\, g^{-1}, \qquad
g \in \Sigma'_j,\ g_1 \text{ son image dans } \mathcal{G}',
\]
formule de la page 10, qu'on vérifie aussitôt : $g\,\Psi_j(\tilde\jmath)(g_1^{-1}
\sigma g_1)\,g^{-1}$ relève $\sigma$ et fixe $g\tilde\jmath$, c'est donc
$\Psi_j(g\tilde\jmath)(\sigma)$.

\pagerange{8}{10}
\section*{« Récapitulons » (pages 8 à 10)}

La page 8 range tout ce qui précède en trois groupes, et la numérotation des
formules repart à (1).

\textbf{I. Les données.} $X$ de type fini sur $\overline{\mathbb{Q}}$ ;
$D = \sum_{i \in I} D_i$ diviseur effectif à composantes simples, $X$ lisse aux
points maximaux de $\operatorname{supp} D$ ; $J \subset X'(\overline{\mathbb{Q}})$,
$X' = X \setminus D$\footnote{« Lisse » ici, « régulier » page 4 : pour un
schéma de type fini sur un corps de caractéristique nulle, c'est la même
chose.}.

\textbf{II. Ce qui s'en déduit.} $\mathcal{G}'$, le groupe des automorphismes
du schéma absolu $X$ qui préservent $D$ et $J$ ; $G' \subset \mathcal{G}'$, ceux
qui sont $\overline{\mathbb{Q}}$-linéaires ; $T$ ; $\Gamma$. D'où
\[
(1)\qquad 1 \to G' \to \mathcal{G}' \to \Gamma, \qquad
(2)\qquad \text{un germe de relèvement } \Gamma' \to \mathcal{G}'.
\]

\textbf{III. Un choix} : un revêtement universel $\widetilde X'$ de $X'$, d'où
$\pi_1(X') = \mathrm{Aut}(\widetilde X'/X')$. On trouve alors :
\begin{enumerate}
\item[(3)] la structure d'extension $1 \to \pi_1(X') \to \Sigma' \to
\mathcal{G}' \to 1$ ;
\item[(4)] les $\Sigma'$-ensembles $\widetilde X'(I)$ et $\widetilde X'(J)$ ;
\item[(5)] les actions de $\Sigma'$ sur $I$ et $J$ à travers $\mathcal{G}'$, et
les $\Sigma'$-applications $\widetilde X'(I) \to I$ et $\widetilde X'(J) \to J
\simeq \widetilde X'(J)/\pi_1$ ;
\item[(6)] l'application $\Phi : \widetilde X'(I) \to \mathrm{Hom}(T,
\pi_1(X'))$, compatible avec $\Sigma'$ ;
\item[(7)] le système $\Psi = (\Psi_j)_{j \in J}$ ;
\item[(9)] un germe d'homomorphisme $\Gamma' \to \Sigma'$ relevant
l'inclusion.
\end{enumerate}
L'élément (8), une variante « extérieure » de (7), est biffé. Deux remarques.
Le germe (9) n'est pas canonique : c'est le $\varphi_\xi$ de (15), et il
dépend du point $\xi$ choisi — pour $\xi \in J$, il se déduit de (2) et de
(7)\footnote{La note encadrée de la page 10 rattache (8) à (7) (« (8) …
“relevant” (7) … ») ; la page n'indique pas de quel point (9) provient. Deux
points distincts donnent des germes non conjugués, si $X'$ se plonge dans une
variété semi-abélienne (page 4).}. Et la marge de (6) prévient : « si $\dim X
= 1$ ; sinon il faut revoir avec plus de soin la situation… » — c'est la
réserve qu'on a faite à la fin de la page 7.

\pagerange{11}{11}
\section*{Le programme des « Gleichnis » (page 11)}

Programme des « gruppentheoretische Gleichnis » — paradigmes en termes de
théorie des groupes — de l'arithmétique\footnote{« Gleichnis » est en allemand
et entre guillemets sur la page ; « gruppentheoretische », ajouté au-dessus de
la ligne, se lit mal. Le mot signifie « parabole », « image », « similitude ».
Une première rédaction biffée disait « de la théorie géométrique algébrique
absolue (en car. 0) ».} :

\begin{enumerate}
\item[a)] donner une description « gruppentheoretisch » complète de $\Gamma$ ;
\item[b)] donner une description complète de tous les éléments de structure
énumérés ci-dessus, dans des cas typiques de $(X, D, J)$, « dont nous
indiquons quelques-uns plus bas ».
\end{enumerate}

Les cas typiques ne viennent pas, sinon, peut-être, sous la forme des espaces
de modules des pages 26 à 32. Une note marginale propose de prendre $X$
projectif pour simplifier.

\emph{NB.} La structure donnée pour $(X, D, J)$ donne automatiquement la même
pour tout revêtement fini étale de $X'$ ; il faudrait l'expliciter avec soin.
Le reste de la note, d'une lecture incertaine, nomme deux opérations : rajouter
certaines des composantes $D_i$ qu'on avait enlevées (« bouclage de
tresses »), ou, si $\dim X = 1$, enlever certains des points de $J$
(« fixage de tresses »)\footnote{Les deux expressions entre guillemets, de lui,
sont des lectures incertaines, comme « explicites » et « fondements » dans la
même phrase. L'interprétation ci-dessous est la nôtre.}. Ce sont les deux
manières de passer d'une courbe trouée à une autre : boucher un trou, ce qui
quotiente $\pi_1(X')$ par le sous-groupe distingué engendré par l'inertie en
ce trou ; ou percer en un point marqué, ce qui en fait un trou de plus. La
seconde est celle que la fibration de la page 32 réalise sur l'espace des
modules.

C'est, sous le nom de « Gleichnis », le programme qu'il appellera dans
l'\emph{Esquisse d'un programme} (1984) et dans sa lettre à Faltings (1983)
\emph{géométrie anabélienne} : l'arithmétique d'une courbe « anabélienne »
serait entièrement codée par son groupe fondamental arithmétique et les
structures qui l'accompagnent\footnote{Le rapprochement est de cette édition ;
le mot « anabélien » n'est pas sur ces pages.}.

\pagerange{12}{13}
\section*{Descente à un corps de nombres, et bitorseurs (pages 12 et 13)}

« RemB »\footnote{Souligné ; lecture douteuse.}. Soit $\Gamma' = \Gamma_K$
pour un corps de nombres $K$, et une forme de $(X, D, J)$ sur $K$ ; elle donne
un scindage partiel $\Gamma' \to \mathcal{G}'$ et, par image inverse, une
extension
\[
(9)\qquad 1 \longrightarrow \pi_1(X') \longrightarrow \Sigma'_{\Gamma'}
\longrightarrow \Gamma' \longrightarrow 1,
\]
qui est le groupe fondamental de $X'_K$. Le groupe $\Gamma'$ agit sur $X'$,
donc sur la catégorie $\mathrm{Rev}(X')$ des revêtements finis étales de $X'$,
qui s'identifie à celle, $\mathrm{Ens}_f(\pi_1)$, des $\pi_1(X')$-ensembles finis
continus. Pour $\gamma \in \Gamma'$, la fibre $\mathcal{B}_\gamma$ de
$\Sigma'_{\Gamma'}$ au-dessus de $\gamma$ est un $\pi_1(X')$-\emph{bitorseur}, par
multiplication à gauche et à droite, et l'action de $\gamma$ sur
$\mathrm{Rev}(X') \simeq \mathrm{Ens}_f(\pi_1)$ s'écrit
\[
(10)\qquad E \longmapsto \mathcal{B}_\gamma \wedge^{\pi_1} E
\]
— à la convention de signe près sur l'action de $\gamma$ sur les revêtements,
l'autre convention remplaçant $\mathcal{B}_\gamma$ par
$\mathcal{B}_{\gamma^{-1}}$. C'est exactement le Corollaire des pages 11 et 12
du dossier 141 : les auto-équivalences de $\mathrm{Ens}(\pi)$ sont les produits
contractés par les bitorseurs\footnote{Le renvoi est de cette édition.}.

Si le revêtement $Y'$ et son transporté ${}^{\gamma}Y'$ se correspondent,
le transport de structure, qui préserve $D$ et $J$, donne des bijections
canoniques
\[
(11)\qquad Y'(I) \simeq ({}^{\gamma}Y')(I), \qquad Y'(J) = ({}^{\gamma}Y')(J),
\]
où $Y'(I)$ désigne les points du normalisé de $X$ dans $Y'$ au-dessus des
$\eta_i$. La page 13 les analyse. Si $Y'$ correspond au $\pi_1$-ensemble $E$,
on a $Y'(I) \simeq \widetilde X'(I) \wedge^{\pi_1} E$ et $Y'(J) \simeq
\widetilde X'(J) \wedge^{\pi_1} E$, et les bijections (11) proviennent
d'isomorphismes de $\pi_1$-ensembles à droite
\[
\widetilde X'(I) \simeq \widetilde X'(I) \wedge^{\pi_1} \mathcal{B}_\gamma,
\qquad \widetilde X'(J) \simeq \widetilde X'(J) \wedge^{\pi_1}
\mathcal{B}_\gamma,
\]
lesquels reviennent à des applications $\mathcal{B}_\gamma \to
\mathrm{Bij}(\widetilde X'(I))$ (et de même pour $J$), compatibles aux actions
de $\pi_1$ à gauche et à droite. Pour $\gamma$ variable, c'est une application
\[
\Sigma'_{\Gamma'} = \bigcup_{\gamma \in \Gamma'} \mathcal{B}_\gamma
\longrightarrow \mathrm{Aut}_{\mathrm{ens}}(\widetilde X'(I)), \qquad
\text{et de même pour } \widetilde X'(J),
\]
et c'est précisément l'action de $\Sigma'$ déjà considérée. Autrement dit,
les $\Sigma'$-ensembles $\widetilde X'(I)$ et $\widetilde X'(J)$ de la
récapitulation sont la même chose qu'un prolongement cohérent, aux pointes
et aux points marqués de tous les revêtements, de l'action de $\Gamma'$ sur
les revêtements\footnote{Cette dernière phrase résume ; la page s'achève sur
« et celles qu'on cherche sont justement l'opération déduite … de $\Sigma'$ sur
$\widetilde X'(I)$, $\widetilde X'(J)$… ». Sous « Aut », sur la troisième
accolade, un mot surchargé ; l'indice « ens » est lu d'après un souligné.}.

\pagerange{14}{16}
\section*{Le groupoïde fondamental sur les points marqués (pages 14 à 16)}

⑦ Il s'agit de reconstruire le groupoïde fondamental $\Pi_1(X'; J)$ de $X'$,
dont les objets sont les points de $J$, et l'action de $\mathcal{G}'$ sur lui,
à partir de l'extension $\Sigma'$ et de l'action de $\Sigma'$ sur $\widetilde
X'(J) \to J$ seulement. Les scindages (7) et (8) sont superflus, dit une note
marginale : « elles résultent des précédentes »\footnote{Le texte les cite
encore, entre crochets, sous la forme $\Psi^{\mathrm{ext}} : \widetilde X'(J)
\to \mathrm{Splitt}(\Sigma'/\mathcal{G}')/\pi_1$ et $\Psi_j$ ; la page 16 se
demande comment ils s'en déduisent. Voici la réponse, de cette édition :
$\pi_1(X')$ agit simplement transitivement sur $\widetilde X'(j)$, de sorte que
le stabilisateur de $\tilde\jmath$ dans $\Sigma'_j$ s'envoie isomorphiquement
sur $\mathcal{G}'_j$ ; $\Psi_j(\tilde\jmath)$ est l'inverse de cet
isomorphisme.}.

Pour $j \in J$, notons $\widetilde X'[j]$ le revêtement universel pointé en
$j$. Un point $\tilde\xi \in \widetilde X'(j)$ l'identifie à $\widetilde X'$ :
$\psi_{\tilde\xi} : \widetilde X' \xrightarrow{\sim} \widetilde X'[j]$, avec
$\psi_{\lambda\tilde\xi} = \psi_{\tilde\xi} \circ \lambda$ pour $\lambda \in
\pi_1$. Les choix se recollent donc en un isomorphisme canonique
\[
\widetilde X'[j] \xleftarrow{\ \sim\ } \widetilde X'(j) \wedge^{\pi_1}
\widetilde X',
\]
où $\widetilde X'(j)$ est un $\pi_1$-torseur à droite. On pose
\[
\mathrm{Isom}_{\Pi_1}(j, j') = \mathrm{Isom}(\widetilde X'[j], \widetilde
X'[j']) \simeq \mathrm{Isom}_{\pi_1\text{-tors.}}\bigl(\widetilde X'(j),
\widetilde X'(j')\bigr),
\]
la composition des chemins étant celle des isomorphismes de
torseurs\footnote{Avec les foncteurs fibres, un chemin de $j$ à $j'$ est un
isomorphisme du foncteur fibre en $j$ vers celui en $j'$, et le lemme de
Yoneda le change en un isomorphisme de $\widetilde X'[j']$ vers $\widetilde
X'[j]$ : la définition de la page est celle des chemins de $j'$ à $j$. Comme
tout groupoïde est isomorphe à son opposé par $f \mapsto f^{-1}$, rien de ce qui
suit n'en dépend. « $\mathrm{Ob}\,\Pi_1 = J$ » est une lecture incertaine du
premier mot.}. Le groupoïde $\Pi_1(X'; J)$ est donc reconstruit à partir de
$\widetilde X'(J)$ et de l'action libre de $\pi_1$ sur lui : c'est ce qu'on
appelle aujourd'hui le \emph{groupoïde de jauge} (ou d'Ehresmann) du
« fibré principal » $\widetilde X'(J) \to J$, de groupe $\pi_1$\footnote{Le nom
est de cette édition. Le groupoïde fondamental sur un ensemble de points de
base est dans SGA 1 et dans le livre de R. Brown (1968) ; Grothendieck en fera
le langage même de l'\emph{Esquisse}.}.

\emph{L'action de $\mathcal{G}'$.} Sur les objets, c'est l'action donnée sur
$J$. Pour $g \in \mathcal{G}'$, soit $\mathcal{B}_g \subset \Sigma'$ la fibre
au-dessus de $g$, un $\pi_1$-bitorseur. L'action de $\Sigma'$ sur $\widetilde
X'(J)$ donne $\widetilde X'(gj) \simeq \widetilde X'(j) \wedge^{\pi_1}
\mathcal{B}_g$\footnote{Dans le crochet, la page écrit « grâce à l'op. de
$\Sigma'$ sur $\widetilde X'(I)$ » ; c'est $\widetilde X'(J)$ qu'il faut.}, d'où
la flèche
\[
\mathrm{Isom}_{\Pi_1}(j, j') \longrightarrow \mathrm{Isom}_{\Pi_1}(gj, gj'),
\qquad f \longmapsto b \circ f \circ b^{-1} \quad (b \in \mathcal{B}_g),
\]
indépendante de $b$ : remplacer $b$ par $b\lambda$ ne change rien, car $f$,
étant un morphisme de torseurs à droite, commute à l'action à gauche de
$\pi_1$\footnote{La page dit « d'où aussitôt la flèche envisagée » ; la
formule explicite est de cette édition.}.

\pagerange{16}{18}
\section*{Reformulation : extensions et groupoïdes (pages 16 à 18)}

Considérons d'une part les systèmes
\[
\bigl(1 \to \pi \to \Sigma \to G \to 1,\ \ \mathfrak{Z}\bigr),
\]
une extension de groupes et un $\Sigma$-ensemble non vide $\mathfrak{Z}$ sur
lequel $\pi$ agit librement : ils forment une catégorie $C$. D'autre part les
systèmes $(\Pi, \widetilde\Pi, p, G, g)$, où $\Pi$ est un groupoïde connexe non
vide, $p : \widetilde\Pi \to \Pi$ un revêtement universel de $\Pi$, et $g : G
\to \mathrm{Aut}(\Pi)$ une action d'un groupe $G$ : ils forment une catégorie
$C'$. On définit des foncteurs quasi-inverses $\alpha : C \to C'$ et $\beta :
C' \to C$\footnote{La page ne définit pas les morphismes de $C$ et $C'$, et ce
qui suit est vérifié objet par objet : $\beta\alpha$ et $\alpha\beta$ redonnent
l'objet de départ à isomorphisme canonique près.}.

\emph{Définition de $\alpha$.} $\mathrm{Ob}\,\Pi = \mathfrak{Z}/\pi$, qu'on
note $J$, et $\mathrm{Isom}_\Pi(j, j') = \mathrm{Isom}_\pi(\mathfrak{Z}_j,
\mathfrak{Z}_{j'})$, où $\mathfrak{Z}_j$ est l'orbite $j$, un $\pi$-torseur ;
$\Pi$ est connexe puisque ces ensembles sont non vides. Le revêtement universel
$\widetilde\Pi$ a pour objets les éléments de $\mathfrak{Z}$, et un seul
morphisme entre deux objets quelconques ; $p$ envoie l'unique flèche $z \to
z'$ sur l'unique $\pi$-isomorphisme $\mathfrak{Z}_j \to \mathfrak{Z}_{j'}$ qui
envoie $z$ sur $z'$. Le groupoïde $\widetilde\Pi$ est donc le groupoïde
\emph{grossier} (indiscret) sur $\mathfrak{Z}$, simplement connexe, et $p$ est
une fibration discrète : $\widetilde\Pi$ est la catégorie des éléments du
foncteur $j \mapsto \mathfrak{Z}_j$\footnote{La page dit « $\widetilde\Pi$ cat.
discrète », et la note de droite précise « cat. discrète sur $\Pi$ définie par
un foncteur $F : \Pi^{\circ} \to \mathrm{Ens}$, $F(j) = \mathfrak{Z}_j$,
$\widetilde\Pi \simeq \Pi_{/F}$ » : « discrète sur $\Pi$ » est ce qu'on appelle
aujourd'hui une fibration discrète, et c'est exact ; comme catégorie,
$\widetilde\Pi$ n'est pas discrète mais grossière. Avec les flèches définies
ci-dessus, $F$ est covariant ; la page écrit $\Pi^{\circ}$, ce qui revient au
même pour un groupoïde, à $f \mapsto f^{-1}$ près.}. On n'a pas encore utilisé
$\Sigma$. Un élément $\sigma \in \Sigma$ agit sur $\mathfrak{Z}$, donc sur le
groupoïde grossier $\widetilde\Pi$, et sur $\Pi$ par
\[
\bigl(f : \mathfrak{Z}_j \to \mathfrak{Z}_{j'}\bigr) \longmapsto \sigma
\circ f \circ \sigma^{-1} : \mathfrak{Z}_{\sigma j} \to \mathfrak{Z}_{\sigma j'},
\]
qui est encore $\pi$-équivariant parce que $\pi$ est distingué. Les éléments de
$\pi$ agissent sur $\widetilde\Pi$ par automorphismes de revêtement, et
trivialement sur $\Pi$ ; d'où l'action de $G = \Sigma/\pi$ sur $\Pi$ par
passage au quotient. C'est la construction de la page 15 rendue abstraite :
$\Pi_1(X'; J) = \alpha(\pi_1(X'), \Sigma', \mathcal{G}', \widetilde X'(J))$.

\emph{Définition de $\beta$.} $\Sigma$ est le groupe des couples $(g,
\bar g)$ formés de $g \in G$ et d'un automorphisme $\bar g$ de $\widetilde\Pi$
qui le recouvre, $\pi = \mathrm{Aut}(\widetilde\Pi/\Pi)$, $\Sigma \to G$ la
projection, et $\mathfrak{Z} = \mathrm{Ob}\,\widetilde\Pi$ avec l'action
évidente. La suite $1 \to \pi \to \Sigma \to G \to 1$ est exacte parce que tout
automorphisme de $\Pi$ se relève au revêtement universel, et $\pi$ agit
librement sur $\mathfrak{Z}$\footnote{Ces deux vérifications sont de cette
édition. Pour $\beta\alpha$ : les automorphismes de la fibration discrète
$\widetilde\Pi \to \Pi$ sont les bijections de $\mathfrak{Z}_j$ qui commutent à
tous les $\pi$-automorphismes de $\mathfrak{Z}_j$, c'est-à-dire l'action de
$\pi$ elle-même.}.

« Variantes profinies… » : le cas des pages précédentes est profini, et la
page s'en tient à cette mention ; l'énoncé vaut pour un groupe profini $\pi$,
des $\pi$-ensembles profinis à action continue et des groupoïdes dont les
ensembles de flèches sont profinis, mais le dossier ne l'écrit pas.

\pagerange{18}{26}
\section*{Les pointes en termes de groupoïdes (pages 18 à 26)}

\pagerange{18}{20}
\subsection*{Le triple $(\mathcal{T}, T, \Phi)$}

⑧ Avec les notations de ⑦, la page se propose d'interpréter de même la donnée
d'un triple $(\mathcal{T}, T, \Phi)$ : $\mathcal{T}$ un $\Sigma$-ensemble, $T$
un groupe — « $\simeq \widehat{\mathbb{Z}}$ ? » — muni d'une action $\chi : G
\to \mathrm{Aut}(T)$, et
\[
\Phi : \mathcal{T} \longrightarrow \mathrm{Hom}_{\mathrm{gr}}(T, \pi)
\]
compatible avec $\Sigma$, qui agit sur le but par $\chi$ et par conjugaison sur
$\pi$. C'est l'abstraction de $\widetilde X'(I)$ et de l'application (6) de la
page 10.

Une première tentative, que deux longs traits obliques traversent de la page
19 au haut de la page 20, considère le quotient $I = \mathcal{T}/\pi$ et le
groupoïde $\langle T, I\rangle$ d'objets $I$ où chaque objet a $T$ pour groupe
d'automorphismes et où deux objets distincts ne sont pas liés, et cherche un
morphisme de ce groupoïde vers le groupoïde des revêtements universels de
$\Pi$, identifié à celui des $\pi$-torseurs. Elle est abandonnée au milieu
d'une phrase\footnote{« … le $\pi$-torseur correspondant $P_j$ s'envoie dans
$\mathcal{T}_j$ …, la fibre de $P_j$ en $t \in \mathcal{T}_j$ étant le torseur
sous le stabilisateur $\pi_t$ … : on considère $\Phi(t) : T \to \pi$ dont
$\pi_t$ centralise ». Les deux traits semblent biffer le passage, qui reste
lisible. La page écrit $J$ pour $\mathcal{T}/\pi$ ; voir les conventions.}.

La seconde tentative restreint à $\pi$ l'action de $\Sigma$ sur $\mathcal{T}$
et forme le \emph{groupoïde d'action} $\langle \mathcal{T}, \pi\rangle$ : ses
objets sont les éléments de $\mathcal{T}$, et $\mathrm{Isom}(t, t') =
\pi_{t,t'} = \{\lambda \in \pi \mid \lambda t = t'\}$. On a
\[
\pi_0\langle \mathcal{T}, \pi\rangle = \mathcal{T}/\pi = I, \qquad
\pi_1(\langle \mathcal{T}, \pi\rangle, t) = \pi_t,
\]
le stabilisateur de $t$ dans $\pi$. Et $\pi_t$ centralise $\Phi(t)(T)$ : pour
$\lambda \in \pi_t$, la compatibilité de $\Phi$ donne $\Phi(t) = \Phi(\lambda
t) = \mathrm{int}(\lambda) \circ \Phi(t)$, $\lambda$ agissant trivialement sur
$T$ puisqu'il est d'image $1$ dans $G$\footnote{La page écrit « on trouve que
$\pi_t$ centralise $\Phi(t)(T)$ » en biffant « centralise », puis un mot
illisible ; l'énoncé est vrai, et la démonstration est de cette édition.}. On
suppose alors que $\Phi(t)$ est un isomorphisme
\[
\Phi_t : T \xrightarrow{\ \sim\ } \pi_t .
\]
Cela force $T$ commutatif — $\pi_t$ centralise alors $\pi_t$ —, ce que la page
suppose à part (« supposons $T$ commutatif »). Dans le cas géométrique, c'est
la propriété des pointes d'une courbe : si $X'$ est une courbe lisse de genre
$g$ privée de $\nu \geq 1$ points, hors du cas $(g, \nu) = (0, 1)$, le
stabilisateur d'une pointe dans $\pi_1(X')$ est son groupe d'inertie, isomorphe
à $T$ — le corps résiduel est $\overline{\mathbb{Q}}$, la décomposition se
réduit à l'inertie — ; en dimension supérieure, il est plus gros (page 7).

\pagerange{21}{24}
\subsection*{$T$-groupoïdes et pseudo-partitions (pages 21 à 24)}

Les isomorphismes $\Phi_t$ forment un système transitif d'identifications de
tous les $\pi_t$ à $T$ : pour $\lambda \in \pi_{t,t'}$, la conjugaison par
$\lambda$ envoie $\pi_t$ sur $\pi_{t'}$, et $\Phi_{t'} = \mathrm{int}(\lambda)
\circ \Phi_t$ ; comme $T$ est commutatif, l'identification ne dépend pas du
chemin. Le groupoïde $\langle \mathcal{T}, \pi\rangle$ devient ainsi un
\emph{$T$-groupoïde} — un groupoïde, pas nécessairement connexe, dont tous les
groupes d'automorphismes sont identifiés à $T$ de manière compatible. On dit
aujourd'hui une \emph{gerbe liée par} le groupe abélien $T$ au-dessus de
l'ensemble discret $I$ ; ici, à équivalence près, une réunion disjointe de
copies du groupe $T$ vu comme groupoïde à un objet\footnote{Le nom (Giraud,
1971) est de cette édition. La page 21, d'une écriture rapide, a beaucoup de
mots illisibles ; elle dit d'abord que $\Sigma$ opère sur ce $T$-groupoïde « en
opérant trivialement sur $T$ », puis « enfin sur $T$ via $G$ ». C'est la
seconde qui est juste : $\Phi_{\sigma t} = \mathrm{int}(\sigma) \circ \Phi_t
\circ \chi(\sigma)^{-1}$, et $\Sigma$ agit sur $T$ par $\chi$ ; l'action n'est
triviale qu'au-dessus du noyau de $\chi$.}. Le groupe $\Sigma$ opère sur ce
$T$-groupoïde, par son action sur $\mathcal{T}$, par conjugaison sur $\pi$ et
par $\chi$ sur $T$ ; et l'on a un foncteur fidèle
\[
\langle \mathcal{T}, \pi\rangle \longrightarrow \langle \pi\rangle, \qquad t
\longmapsto \ast, \quad \lambda \in \pi_{t,t'} \longmapsto \lambda,
\]
vers le groupoïde à un objet de groupe $\pi$, qui est, à équivalence près, le
groupoïde des revêtements universels de $\Pi$.

\emph{Réciproquement} (pages 23 et 24), donnons-nous un $T$-groupoïde
$\Lambda$, une action de $\Sigma$ sur $(T, \Lambda)$ et un foncteur fidèle
$\Lambda \to \langle \pi\rangle$ compatible avec $\Sigma$ ; posons $\mathcal{T}
= \mathrm{Ob}\,\Lambda$, et notons $\pi_{t,t'} \subset \pi$ l'image de
$\mathrm{Isom}_\Lambda(t, t')$. Pour retrouver une action de $\pi$ sur
$\mathcal{T}$, la page exige
\[
\forall t \in \mathcal{T}, \quad \text{la famille } (\pi_{t,t'})_{t' \in
\mathcal{T}} \text{ est une pseudo-partition de } \pi,
\]
c'est-à-dire que tout $\lambda \in \pi$ est dans un et un seul $\pi_{t,t'}$ :
il y a une et une seule flèche de $\Lambda$ d'origine $t$ au-dessus de
$\lambda$. C'est la condition de relèvement unique des chemins — la marge dit
« condition $\exists!$ de remontée » — qui fait de $\Lambda \to \langle\pi\rangle$
un \emph{revêtement de groupoïdes}. On pose alors $\lambda t$ = le but de ce
relèvement, et l'on vérifie $\Lambda \simeq \langle \mathcal{T}, \pi\rangle$ :
c'est l'équivalence classique entre revêtements du groupoïde
$\langle\pi\rangle$ et $\pi$-ensembles\footnote{Le nom est de cette édition ;
voir R. Brown, \emph{Topology and Groupoids}. « Pseudo-partition » est une
lecture incertaine : une partition dont des membres peuvent être vides.}. Si
de plus les stabilisateurs $\pi_t$ sont abéliens, on retrouve une famille de
groupes indexée par $I = \mathcal{T}/\pi$, et l'action de $\Sigma$.

« Ce n'est pas bien éclairant. »

\pagerange{24}{26}
\subsection*{Un foncteur sur les revêtements, et l'arrêt (pages 24 à 26)}

« … peut-être est-il plus agréable » de considérer $\mathcal{T}$ comme un
foncteur. Les revêtements de $\Pi$ forment une catégorie
$\mathrm{Rev}(\Pi)$ équivalente à celle des $\pi$-ensembles, par $E \mapsto
\widetilde\Pi^{0} \times_\pi E$, et l'on considère
\[
F : (\pi\text{-}\mathrm{Ens}) \simeq \mathrm{Rev}(\Pi) \longrightarrow
(\mathrm{Ens}), \qquad E \longmapsto \mathcal{T}^{0} \wedge_\pi E,
\]
où l'indice ${}^{0}$ marque, semble-t-il, que l'action à gauche de $\pi$ est
changée en action à droite\footnote{La note qui explique l'indice est presque
illisible ; l'interprétation est de cette édition, et c'est celle que demande
le produit contracté.}. C'est un foncteur qui commute aux limites inductives
quelconques, et tout tel foncteur est de cette forme : c'est la proposition des
pages 11 et 12 du dossier 141, où un foncteur $\mathrm{Ens}(\pi) \to
\mathcal{C}$ commutant aux limites inductives équivaut à un $\pi$-objet à droite
de $\mathcal{C}$. On veut ensuite que $F$ soit compatible avec l'action de
$\Sigma$ sur $\mathrm{Rev}(\Pi)$, c'est-à-dire des isomorphismes $F({}^{\gamma}
\Pi') \simeq F(\Pi')$, fonctoriels en $\Pi'$ et transitifs en $\gamma$ ; par le
calcul de la page 13 (les bitorseurs $\mathcal{B}_\gamma$), cela revient à
prolonger l'action de $\pi$ sur $\mathcal{T}$ en une action de $\Sigma$.

« Mais ceci ne donne pas d'interprétation … de la $T$-[structure]. Arrêtons
… » : l'application $\Phi$, qui est la vraie donnée kummérienne, reste hors
de cette formulation, et le dossier ne la reprend pas\footnote{Le mot qui suit
« $T$- » est illisible.}.

\pagerange{26}{27}
\section*{Topos modulaires de genre $0$ (pages 26 et 27)}

« ⑨ Les schémas et Topos fondamentaux dans la théorie… », et, en
sous-titre, « (A) Topos modulaires ». Le topos modulaire
$\mathcal{T}_{0,\nu}$ — on dit aujourd'hui le champ $M_{0,\nu}$ — classe les
courbes de genre $0$ munies de $\nu$ points marqués distincts.

\emph{Pour $\nu = 0, 1, 2$}, c'est le topos classifiant d'un groupe algébrique :
\[
M_{0,0} = B\,\mathrm{PGL}(2), \qquad M_{0,1} = B\,\mathrm{Aff}(1), \qquad
M_{0,2} = B\,\mathbb{G}_m,
\]
parce que toutes les courbes de genre $0$ à $\nu$ points marqués sont
localement isomorphes à $(\mathbb{P}^1; \infty)$, $(\mathbb{P}^1; 0, \infty)$,
etc., et que le groupe de leurs automorphismes est de dimension $3 - \nu > 0$.
Ce ne sont donc pas des champs de Deligne--Mumford — ce que la page 28 dit
en les excluant\footnote{La page 26 est lacunaire : « … pour $\nu = 0, 1, 2$
existe [comme topos] étendues — car les schémas … d'hom. d'une $X/S$ fibré
projectif de genre $0$ et de $\nu$ sections marquées … Cela tient au fait que …
$3 - \nu > 0$ ». La transcription ne permet pas de décider s'il affirme ou s'il
nie qu'ils soient des « étendues » ; ce qui est vrai est que ce sont des champs
algébriques d'Artin à stabilisateurs de dimension positive, et la page 28 les
range parmi les exceptions. Pour $\nu = 0$, une courbe de genre $0$ n'est que
localement (pour la topologie étale) isomorphe à $\mathbb{P}^1$ : c'est une
conique.}. Leurs groupes fondamentaux géométriques, comme celui de $M_{0,3} =
\mathrm{Spec}\,\mathbb{Q}$, sont triviaux :
\[
\pi_1(\overline{M}_{0,\nu}) = 1, \qquad 0 \leq \nu \leq 3,
\]
le groupe fondamental du classifiant d'un groupe connexe étant
trivial\footnote{La page écrit $\pi_1(\mathcal{T}_{0,\nu}) = 0$, le $3$ repassé
sur un autre chiffre ; c'est le groupe fondamental sur $\overline{\mathbb{Q}}$
qui est trivial. Sur $\mathbb{Q}$, c'est $\Gamma$.}.

\emph{Pour $\nu \geq 4$}, en envoyant les trois premiers points sur $0, 1,
\infty$, $M_{0,\nu}$ devient le schéma des systèmes $(x_1, \ldots, x_{\nu-3})$
de $\nu - 3$ points distincts de $\mathbb{P}^1 \setminus \{0, 1, \infty\}$.
Son groupe fondamental géométrique est le complété profini d'un groupe de
tresses $\Pi_{0,\nu}$ : précisément, du groupe modulaire pur de la sphère à
$\nu$ trous, qui est le groupe des tresses pures à $\nu - 1$ brins quotienté
par son centre\footnote{La page dit « un groupe de tresses », lecture incertaine,
sans préciser ; la précision est de cette édition. Le passage du groupe
discret au complété est le théorème d'existence de Riemann.}. Il dépend
contravariamment de $[1, \nu]$ pour les injections : une injection $[1, \nu']
\hookrightarrow [1, \nu]$ définit l'oubli $M_{0,\nu} \to M_{0,\nu'}$ des
autres points, et le morphisme induit sur les $\pi_1$. Le groupe
$\mathfrak{S}_\nu$ agit sur $\overline{M}_{0,\nu}$ en permutant les points, par
des automorphismes définis sur $\mathbb{Q}$ ; donc $\mathfrak{S}_\nu \times
\Gamma$ agit sur $\overline{M}_{0,\nu}$ comme schéma absolu. La page s'arrête
là, au tiers.

\pagerange{28}{31}
\section*{Une conjecture arithmético-géométrique (pages 28 à 31)}

\pagerange{28}{29}
\subsection*{La conjecture (pages 28 et 29)}

Soient $g, \nu \geq 0$ et $M_{g,\nu}$ le topos modulaire sur $\mathbb{Q}$ des
courbes lisses propres, géométriquement connexes, de genre $g$, munies de $\nu$
sections ordonnées deux à deux disjointes. Sauf pour $(g, \nu) = (0, 0), (0,
1), (0, 2), (1, 0)$, c'est une « étendue algébrique » — un champ de
Deligne--Mumford : ce sont exactement les cas $2g - 2 + \nu > 0$, où les
automorphismes d'une courbe marquée sont en nombre fini\footnote{La page écrit
d'abord une autre lettre, qui semble un $\mathcal{T}$, puis $M$ par-dessus ;
« associé » et « étendue algébrique » sont des lectures incertaines. Le nom de
champ de Deligne--Mumford (1969) est de cette édition.}. On suppose désormais
$2g - 2 + \nu > 0$.

Soit $\xi$ un point géométrique de $M_{g,\nu}$, $(X; x_1, \ldots, x_\nu)$ la
courbe marquée sur $\overline{\mathbb{Q}}$ correspondante, et
$\widehat\pi_{g,\nu} = \pi_1(X \setminus \{x_1, \ldots, x_\nu\})$. La
courbe universelle, privée de ses $\nu$ sections, est une fibration au-dessus
de $M_{g,\nu}$ de fibre $X \setminus \{x_i\}$ en $\xi$ ; l'action de monodromie
donne des homomorphismes naturels
\[
(*)\qquad \pi_1(M_{g,\nu}; \xi) \longrightarrow
\mathrm{Out}^{\circ}_{\mathrm{lac}}(\widehat\pi_{g,\nu}),
\qquad
(**)\qquad \pi_1(\overline{M}_{g,\nu}; \xi) \longrightarrow
\mathrm{Out}^{\circ\circ}_{\mathrm{lac}}(\widehat\pi_{g,\nu}),
\]
où ${}^{\circ}$ restreint aux automorphismes extérieurs qui préservent chacune
des classes d'inertie aux $x_i$ — à multiplicateur près —, et ${}^{\circ\circ}$
à ceux qui induisent de plus l'identité sur $T$\footnote{La page dit « qui …
les classes d'inertie aux $x_i$ » et, pour ${}^{\circ\circ}$, « … strictement
… induisant l'identité sur $T = \mathrm{Tate}(\overline{\mathbb{Q}})$ », avec
plusieurs mots illisibles ; les définitions ci-dessus sont celles que le reste
du dossier impose.}. C'est, pour la courbe marquée universelle, la flèche (19)
de la page 6. En théorie topologique, l'analogue
\[
(***)\qquad \pi_1^{\mathrm{orb}}\bigl(M_{g,\nu}(\mathbb{C})\bigr)
\longrightarrow \mathrm{Out}^{\circ\circ}_{\mathrm{lac}}\bigl(\pi_1(X \setminus
\{x_1, \ldots, x_\nu\})\bigr) = T_{g,\nu}
\]
est un isomorphisme : l'espace de Teichmüller étant contractile, le premier
groupe est le groupe modulaire pur $T_{g,\nu}$, et le théorème de
Dehn--Nielsen--Baer l'identifie au groupe des automorphismes extérieurs du
groupe de surface qui préservent chaque classe de lacet périphérique et
l'orientation\footnote{La page met l'accolade « $= \mathcal{T}_{g,\nu}$
(Teichmüller) » sous $\pi_1(X \setminus \{x_1, \ldots, x_\nu\})$ ; c'est le
groupe des automorphismes extérieurs qui est $T_{g,\nu}$. Dehn--Nielsen--Baer
(années 1920) lui était connu ; il parle de « théorie top. ».}. Par
comparaison, $\pi_1(\overline{M}_{g,\nu}) \simeq \widehat T_{g,\nu}$, d'où le
diagramme (4) :

\begin{tikzcd}[column sep=small, row sep=normal, nodes={font=\scriptsize}]
1 \arrow[r] & \pi_1(\overline{M}_{g,\nu}) \arrow[r] \arrow[d, "\wr"'] &
\pi_1(M_{g,\nu}) \arrow[r] \arrow[d] & \Gamma \arrow[r] \arrow[d, "\chi"] & 1 \\
1 \arrow[r] & \widehat T_{g,\nu} \arrow[r] &
\mathrm{Out}^{\circ}_{\mathrm{lac}}(\widehat\pi_{g,\nu}) \arrow[r] &
\widehat{\mathbb{Z}}^{*} &
\end{tikzcd}

où la flèche de droite de la seconde ligne est le multiplicateur\footnote{La
page laisse le dernier terme de la seconde ligne coupé par le bord du
feuillet (« $\mathrm{Out}\ldots$ ») et trace une flèche oblique depuis $\Gamma$ ;
le multiplicateur est ce que la ligne impose, et la flèche oblique est alors
$\chi$. La flèche de gauche est marquée « $\wr$ ? ».}. La seconde ligne n'est
exacte en son milieu que si $\widehat T_{g,\nu} \to
\mathrm{Out}^{\circ\circ}_{\mathrm{lac}}(\widehat\pi_{g,\nu})$ est surjective, et
à gauche que si elle est injective : la page note « … que $\widehat T_{g,\nu}
\to \mathrm{Out}^{\circ}_{\mathrm{lac}}$ injective … ». Cette injectivité est
ce qu'on appelle aujourd'hui le \emph{problème des sous-groupes de congruence}
pour les groupes modulaires ; elle est établie pour $g \leq 2$, et ouverte en
général à notre connaissance\footnote{À notre connaissance : Asada (2001) pour $g \leq 1$, Boggi pour
$g = 2$ ; les références sont de cette édition.}.

\emph{La conjecture.} « On aimerait » que la flèche verticale du milieu soit un
isomorphisme. Cela donnerait une construction de $\Gamma$ comme
\[
\Gamma \simeq \mathrm{Out}^{\circ}_{\mathrm{lac}}(\widehat\pi_{g,\nu}) \big/
\widehat T_{g,\nu}
\]
pour tout $(g, \nu)$, à partir d'un groupe « bien connu ». \emph{Exemple} : pour
$g = 0$, $\nu = 3$, $\widehat T_{0,3} = 1$, et l'on obtient l'expression « déjà
envisagée »
\[
\Gamma \simeq \mathrm{Out}^{\circ}_{\mathrm{lac}}\bigl(\widehat\pi_1(\mathbb{P}^1
\setminus \{0, 1, \infty\})\bigr),
\]
le groupe des automorphismes extérieurs à lacets purs du groupe libre profini
de rang $2$ — celui qu'étudie le dossier 145, où il est noté
$\mathrm{Out}^{!}_{\mathrm{lac}}(\widehat\pi)$\footnote{Les renvois de la page,
« la flèche (1) (la flèche verticale centrale de (5)) », sont lus tels quels ;
le diagramme est numéroté (4). Une note marginale propose une forme plus
prudente : « Vérifier conjecture : $\pi_1(M_{g,\nu;\xi}) \simeq$ normalisateur de
$\widehat T_{g,\nu}$ dans $\mathrm{Out}^{\circ}_{\mathrm{lac}}$ ». Une autre,
encadrée et raturée, rappelle les exceptions.}. L'injectivité de $\Gamma$
dans ce groupe est le théorème de Belyi (1979). Que l'image soit tout le groupe
n'est pas démontré, et ce n'est pas la forme sous laquelle la question a été
poursuivie : Drinfeld (1990) et Ihara (1991) ont défini le groupe de
Grothendieck--Teichmüller $\widehat{GT}$, intermédiaire, par des relations
issues de $M_{0,4}$ et $M_{0,5}$, et l'égalité $\Gamma = \widehat{GT}$ est
ouverte\footnote{Ce paragraphe est de cette édition ; aucun de ces noms, sauf
Belyi, n'est sur les pages, et Drinfeld comme Ihara sont postérieurs.}.

\pagerange{30}{31}
\subsection*{L'action des permutations, et deux questions (pages 30 et 31)}

Soit $N_{g,\nu} = M_{g,\nu}/\mathfrak{S}_\nu$ le champ des courbes munies de
$\nu$ points non ordonnés ; $M_{g,\nu} \to N_{g,\nu}$ est un revêtement étale
galoisien de groupe $\mathfrak{S}_\nu$, d'où
\[
(5)\qquad 1 \longrightarrow \pi_1(\overline{M}_{g,\nu}) \longrightarrow
\pi_1(N_{g,\nu}; \xi) \longrightarrow \mathfrak{S}_\nu \times \Gamma
\longrightarrow 1,
\]
où le groupe du milieu s'écrit aussi $\pi_1(\overline{M}_{g,\nu},
\mathfrak{S}_\nu \times \Gamma)$ : c'est l'extension $\Sigma$ de la page 2,
pour le schéma (ou le champ) $\overline{M}_{g,\nu}$ et le groupe
$\mathfrak{S}_\nu \times \Gamma$ d'automorphismes absolus de la page
27\footnote{Le rapprochement avec la page 2 est de cette édition ; la page
pose les trois écritures sous le terme du milieu, reliées par des $=$
verticaux.}. La conjecture équivaut à
\[
\pi_1(N_{g,\nu}; \xi) \xrightarrow{\ \sim\ }
\mathrm{Out}_{\mathrm{lac}}(\widehat\pi_{g,\nu}),
\]
où l'on permet maintenant de permuter les classes d'inertie : les deux
groupes sont extensions de $\mathfrak{S}_\nu$ par les deux membres de $(*)$,
et toute permutation des trous est réalisée par un élément du groupe
modulaire. On en tirerait $\Gamma \simeq
\mathrm{Out}_{\mathrm{lac}}(\widehat\pi_{g,\nu})/\widehat{\widetilde T}_{g,\nu}$, où
$\widetilde T_{g,\nu} = \mathrm{Out}^{+}_{\mathrm{lac}}(\pi_1(X^{\mathrm{top}}
\setminus \{x_i\}))$ est le groupe modulaire complet, qui permute les
trous\footnote{L'exposant de $X_\xi$ est lu « top » sous réserve.}.

\emph{NB}, d'une autre encre : il faudrait déterminer le groupe
$G_{g,\nu}$ des automorphismes du topos modulaire géométrique
$\overline{M}_{g,\nu}$ — est-ce l'image de $\mathfrak{S}_\nu$ ? —, sauf que
$\mathfrak{S}_\nu$ n'agit pas fidèlement pour $g = 0$, $\nu \leq 4$. C'est
juste pour $\nu = 4$ : $M_{0,4} = \mathbb{P}^1 \setminus \{0, 1, \infty\}$, et
$\mathfrak{S}_4$ y agit à travers $\mathfrak{S}_3$, le groupe anharmonique des
six transformations du dossier 145, avec pour noyau le groupe de Klein ; et
pour $\nu = 3$, $\overline{M}_{0,3}$ est un point\footnote{L'identification du noyau est
de cette édition.}.

\emph{Question} :
\[
\Gamma \longrightarrow \mathrm{Out}(\widehat T_{g,\nu}) =
\mathrm{Out}\bigl(\pi_1(\overline{M}_{g,\nu})\bigr) \quad \text{est-il
injectif ? bijectif ?}
\]
« Sans doute il faut ajouter [une] condition, cf. NB. » Pour $(g, \nu) = (0,
4)$, $\widehat T_{0,4}$ est le groupe libre profini de rang $2$ : l'injectivité
est le théorème de Belyi, et la bijectivité est fausse, car un automorphisme
extérieur qui ne préserve pas les classes d'inertie — $x \mapsto xy$, $y \mapsto
y$ — n'est pas dans l'image. La condition à ajouter est donc bien du type
« à lacets »\footnote{Exemple de cette édition.}.

\emph{Question} (page 31, seule ligne de la page) : $\mathrm{Out}(T_{g,\nu}) =
1$ ? Pour le groupe discret, la réponse est non : la conjugaison par une
classe renversant l'orientation est un automorphisme extérieur non trivial, et
pour $g \geq 3$, $\nu = 0$, c'est le seul — $\mathrm{Out}(T_g) = \mathbb{Z}/2$
(Ivanov, McCarthy, 1986--1988, après ces pages) ; pour le groupe pur, la
conjugaison par les classes qui permutent les trous en fournit en général
d'autres\footnote{La lecture de la question est sûre ; la réponse est de
cette édition.}.

\pagerange{32}{33}
\section*{La fibration d'oubli d'un point (page 32)}

Sur un listing d'ordinateur retourné, en largeur, deux diagrammes « en
perspective ». Le premier est celui de l'oubli du dernier point, $M_{g,\nu+1}
\to M_{g,\nu}$, qui est la courbe universelle privée de ses $\nu$ sections ;
sa fibre en $\xi$ est $X_\xi \setminus S_\nu$, $S_\nu = \{x_1, \ldots, x_\nu\}$.
Pour $2g - 2 + \nu > 0$, le groupe fondamental de la fibre est de centre
trivial, et l'on a un diagramme à lignes et colonnes exactes :

\begin{tikzcd}[column sep=small, row sep=small, nodes={font=\scriptsize}]
& 1 \arrow[d] & 1 \arrow[d] & & \\
& \pi_1(\overline{X}_\xi \setminus S_\nu) \arrow[r, "="] \arrow[d] &
\pi_1(\overline{X}_\xi \setminus S_\nu) \arrow[d] & & \\
1 \arrow[r] & \pi_1(\overline{M}_{g,\nu+1}) \arrow[r] \arrow[d] &
\pi_1(M_{g,\nu+1}) \arrow[r] \arrow[d] & \Gamma \arrow[r] \arrow[d, "="] & 1 \\
1 \arrow[r] & \pi_1(\overline{M}_{g,\nu}) \arrow[r] \arrow[d] &
\pi_1(M_{g,\nu}) \arrow[r] \arrow[d] & \Gamma \arrow[r] & 1 \\
& 1 & 1 & &
\end{tikzcd}

Sa colonne de gauche est la version profinie de la suite exacte de Birman
(1969)\footnote{La page dessine chaque ligne décalée vers la droite ; le
diagramme est redessiné ici sur une grille. Le nom est de cette édition.}.

Le second diagramme lui fait face du côté des groupes d'automorphismes : la
colonne $1 \to \widehat\pi_{g,\nu} \to \widehat T_{g,\nu+1} \to \widehat
T_{g,\nu} \to 1$, et, au milieu, $1 \to \widehat\pi_{g,\nu} \to
\mathrm{Aut}^{\circ}_{\mathrm{lac}}(\widehat\pi_{g,\nu})^{!} \to
\mathrm{Out}^{\circ}_{\mathrm{lac}}(\widehat\pi_{g,\nu})^{!} \to 1$, les automorphismes
intérieurs formant le noyau parce que $\widehat\pi_{g,\nu}$ est de centre
trivial ; les deux lignes horizontales se terminent sur $\Gamma$, reliées par
un double trait, et portent loin à droite les étiquettes $\Sigma_{\nu+1}$ et
$\Sigma_\nu$\footnote{Sous la flèche vers $\Gamma$, un mot souligné d'une
accolade se lit, mal, « tresses ». L'exposant $!$ n'est pas défini ; au dossier
145, il marque la partie « pure ». Les étiquettes $\Sigma_{\nu+1}$, $\Sigma_\nu$
désignent vraisemblablement les extensions $\pi_1(M_{g,\nu+1})$,
$\pi_1(M_{g,\nu})$ de la ligne correspondante, dans la notation $\Sigma$ des
pages 1 et 2 ; ce n'est qu'une lecture.}. Ce que le dessin exprime, et qui est
vrai : la fibre ayant un point base naturel — le $(\nu + 1)$-ième point —, la
monodromie de $M_{g,\nu+1}$ tombe dans $\mathrm{Aut}$, non seulement dans
$\mathrm{Out}$, et $\pi_1(M_{g,\nu+1})$ est le produit fibré de
$\pi_1(M_{g,\nu})$ et de $\mathrm{Aut}(\widehat\pi_{g,\nu})$ au-dessus de
$\mathrm{Out}(\widehat\pi_{g,\nu})$ : c'est la conséquence formelle de la
trivialité du centre.

C'est aussi, sans que la page le dise, ce qui rattache la fin du dossier à son
début. Un point $\xi \in M_{g,\nu}(K)$ donne un germe de section $\Gamma_K \to
\pi_1(M_{g,\nu})$, le $\varphi_\xi$ de la page 3 ; l'image inverse par ce germe
de la ligne du milieu est l'extension $\Sigma'_{\Gamma'}$ de la page 12 pour la
courbe $X' = X_\xi \setminus S_\nu$ ; et ses $\nu$ sections fournissent les
pointes de la page 6. Les données du paradigme d'une courbe sont celles de la
courbe universelle, ramenées en un point\footnote{Ce paragraphe est de cette
édition.}.

Au bas du feuillet, deux fragments : $M'_{g,\nu}(k) \ni \xi$,
$\mathcal{T}'_{g,\nu} \ni G_\xi$, et $\widehat{\widetilde\Sigma}_{g,\nu} =
\overline{\Sigma}_{g,\nu}$ à côté de $\Sigma_{g,\nu}$, précédés de signes
raturés ; on n'en tire rien. La page 33 est un listing de nombres, sans
écriture de lui sinon deux sommes numériques étrangères au sujet.

\end{document}
